\chapter{Einleitung}
\label{chap:einleitung}

Dieses Dokument dient als Leitfaden für die Anfertigung wissenschaftlicher Arbeiten am Institut für Automatisierungstechnik. Es ist in der Form verfasst, in der die wissenschaftlichen Arbeiten verfasst werden sollen. Die Lücken sind mit den notwendigen Informationen zu füllen, bzw. die Erklärungen zu ersetzen.
Grundsätzlich gilt neben diesem Dokument die aktuelle Fassung Ihrer Prüfungsordnung. Die Vorgaben aus der Prüfungsordnung haben Vorrang, sofern sie im Widerspruch zu diesem Dokument stehen. In diesem Fall informieren Sie bitte Ihren Betreuer hierüber.
Es gibt eine Vielzahl an Literatur die sich mit der Form und Durchführung der wissenschaftlichen Arbeit beschäftigt, auf diese sei hier kurz verwiesen. Hierin werden alle Probleme und Fragestellungen die im Laufe einer Arbeit auftreten können ausführlich behandelt.
Der Abschnitt \ref{chap:allgemeines} befasst sich mit allgemeinen Vorgaben des Instituts. In Abschnitt 3 folgen dann Vorga-ben an die Form und das Layout der Arbeit. Abgeschlossen wird dieses Dokument von weiteren Er-gänzungen im Abschnitt 4. Das Literaturverzeichnis und der Anhang sind nur der Vollständigkeit halber angefügt. 